% The Memorandum of Mere Catholicity, published as its own paper since
% version 1.1. Same source conventions as confession.tex: reproducible
% build, anonymity hygiene, house palette.
\documentclass[11pt,letterpaper]{article}
\ifdefined\pdfsuppressptexinfo\pdfsuppressptexinfo=-1 \fi
\ifdefined\pdftrailerid\pdftrailerid{}\fi
\usepackage[margin=1in]{geometry}
\usepackage[T1]{fontenc}
\usepackage[utf8]{inputenc}
\usepackage{mathpazo}
\usepackage{microtype}
\usepackage{xcolor}
\usepackage{needspace}
\usepackage[colorlinks=true,linkcolor=ink,urlcolor=maroon,citecolor=ink,filecolor=ink]{hyperref}
\urlstyle{same}
\hypersetup{pdftitle={Memorandum},pdfauthor={Anonymous}}
\linespread{1.06}
\setlength{\emergencystretch}{1.5em}

\definecolor{ink}{HTML}{2A2521}
\definecolor{heading}{HTML}{8C4A32}
\definecolor{subhead}{HTML}{6E5642}
\definecolor{accent}{HTML}{B0894C}
\definecolor{maroon}{HTML}{7B2E2E}

\newcommand{\unit}[1]{\needspace{3\baselineskip}\par\medskip{\normalfont\bfseries\color{subhead}#1}\par\nobreak\smallskip}

\begin{document}
\color{ink}

\begin{center}
{\LARGE\bfseries\color{heading} Memorandum\par}
\vspace{0.8em}
{\itshape\color{subhead} An open letter on the reception of the disputed
definitions, with a proposal. To be read before the paper it prefaces.\par}
\vspace{0.5em}
{\color{accent}\hrule height 1pt}
\end{center}
\vspace{1em}

\unit{To.}
The Dicastery for Promoting Christian Unity. The Ecumenical Patriarchate and the holy synods of the Orthodox churches. The synods of the Oriental Orthodox churches. The Assyrian Church of the East. The Saint Irenaeus Joint Working Group, the North American Orthodox-Catholic Consultation, and the standing dialogue commissions. The address matches the argument, for the paper that follows holds that all four ancient communions carry the deposit, so all four are addressed.

\unit{Communicated also.}
To the churches of the Union of Scranton, whose declaration of 2008 professes nearly what this paper professes, and whose part in the act proposed is therefore not to profess but to be received back, the proposal's first test case. To the Anglicans of the high and orthodox parishes, whose contested orders this paper grades an open question, healable by the party that raised the doubt, as Rome healed one in 2001. And in brotherhood to the confessional Lutheran and Reformed churches, who hold real planks of this confession and whose witness belongs to the \emph{sensus fidelium} of the whole Body. The mutual act is not theirs to enact, for the professions proposed belong to the synods of the ancient communions. But the paper names these churches, grades their standing, and owes them the courtesy of address rather than report.

\unit{The thesis.}
A council binds when the whole Body receives it, and reception is material, the coming-to-confess of the faith defined, not the formal acceptance of a formula, an assembly, or an act. On that definition the record shows something the standard account misses. The Body has been torn since 251, at Ephesus and Chalcedon it tore again, yet the seven councils stand, and reception across the open tears is attested wherever it has been checked. The church that refused Ephesus confesses, in Babai its doctor and by common declaration of 1994, the one Christ that Ephesus defined. The churches that refused Chalcedon confess, on first-millennium evidence and by the agreed statements since, the faith it guarded. Binding by act alone never existed in the conciliar age. Every council's authority completed by reception across whatever tears then stood. What division changes is the speed and legibility of the verdict, not its mechanism. It follows that the definitions each communion has made since the great tears, Vatican I and the Marian definitions on the Roman side, the Palamite councils on the Orthodox, stand not yet received by the whole, and not yet rejected, and that no single communion, ours included, can close those questions for the whole by any act of its own.

\unit{The proposal.}
The instrument exists and the Melkite synod of 1995 supplied its genre, a synodal profession, here made mutual and multilateral. Each communion, by the synod competent in it, professes three things. The shared deposit of the councils. That its post-schism definitions are held as its own theological patrimony, taught in its schools and prayed in its liturgies, not imposed on the other communions as conditions of communion. And recognition of the others' baptism, orders, and eucharist, in force from Rome's side fully in law, and from the East's in much though divided practice. Between the Eastern Orthodox and the Oriental churches that recognition is itself the live work of the third profession, so the act stages rather than presumes a single day. Communion follows the professions. No joint definition is asked, for the professions define nothing. They receive. The contested definitions then stand before the reunited Body as proposals of the churches that made them, to ripen or wait as confession spreads or does not. Reception of this proposal is visible and falsifiable, communion restored see by see, names restored to the diptychs, the professions echoed in catechisms across a generation. So is rejection, standing condemnation renewed rather than lapsed.

\unit{What this asks, and its price.}
No church is asked to renounce the content of what it teaches. Rome is not asked to unsay what 1870 defined, nor the Orthodox what 1351 defined. Each is asked to cease requiring of the whole what the whole has not confessed, which the paper argues is all the first millennium itself ever required. For Rome the asking is a real reversal of practice and of self-understanding, for Pastor Aeternus claims a binding force owing nothing to the Church's consent, and to hold that definition as patrimony rather than exact it as a condition is to walk part of the way back toward the first millennium. Section 4 prices this collision. It argues only that the walk is of a kind Rome has already made where the second millennium overreached, at Florence's rigor on the unbaptized, at the coercion of conscience, at the ban on praying with the separated, and made well. What is asked is not the surrender of a possession but the retirement of a claim to one. And the retirement returns more than it costs. A see that may not be corrected spends her credibility with every correction she in fact makes. A see whose acts await the whole Body's reception can correct herself and remain herself. The Orthodox are asked the mirror price, that the identification of the whole Church with themselves be held no longer as a condition. The end sought is not architectural tidiness. It is the Lord's own prayer, that they may all be one, as the Father and the Son are one, so that the world may believe (John 17:21). The paper is an argument that this road is the one the record leaves open.

\unit{The standing invitation.}
John Paul II asked in Ut Unum Sint, sections 95 and 96, for help in finding a way of exercising the primacy open to a new situation, and The Bishop of Rome of 2024 gathered the answers to that invitation. The proposal above is tendered as part of that solicited contribution, for it concerns the exercise. The invitation does not cover the definitions themselves, and we do not pretend it does. Where the invitation runs out, the next paragraph begins.

\unit{The costs, self-assessed.}
The authors are Roman Catholics. By Rome's present law this position is public material dissent, and we do not claim she permits it. We hold it as an appeal to a reception the whole Body has not yet given, and we stand at her altar bearing what follows. The paper's weighings, that Vatican I differs in kind from the primacy the first millennium granted, that the Orthodox exclusivist claim rests on a boundary tradition the undivided Church declined in its strict form, are made by fallible readers and are submitted into the same reception they describe. The paper grades its own evidence, states its open questions in a closing list, and ends with appendices, one grading every act weighed with where the weighing is argued, one applying the rules to the hard cases.

\unit{The ask.}
That the dialogue bodies examine the material-reception thesis and the proposal of mutual professions, and that the paper be read as what it is, a working paper in eight sections written for this examination, opening with the confession it defends and closing with appendices that grade every act weighed and apply the rules to the hard cases, and, after the appendices, with an afterword of counsel for the layman, which the commissions may pass over and the persuaded reader may not.

\vspace{2em}
{\small\noindent This work is dedicated to the public domain (Creative
Commons CC0 1.0). No rights reserved. It may be copied, translated,
printed, and distributed freely, by anyone, for any purpose.
\href{https://merecatholicity.com}{merecatholicity.com}\par}

\end{document}
